\documentclass{beamer}
\usetheme[faculty=econ]{fibeamer}

\usepackage[utf8]{inputenc}
\usepackage[francais]{babel}
\usepackage[T1]{fontenc}


\usepackage{ulem}
\usepackage{listingsutf8}

\usepackage{lstlangarm}
\usepackage{xcolor}
% For multirow tables
\usepackage{multirow}
\definecolor{CommentGreen}{rgb}{0,.6,0}
\lstdefinestyle{arm}
{
  language=[arm]Assembler,
  escapechar=@@, % include LaTeX code between `@' characters
  keepspaces,   % needed to preserve spacing with lstinline
  basicstyle=\scriptsize\ttfamily\bfseries,
  commentstyle=\color{CommentGreen},
  stringstyle=\color{cyan},
  showstringspaces=false,
  keywordstyle=[1]\color{blue},    % instructions
  keywordstyle=[2]\color{magenta}, % directives
  keywordstyle=[3]\color{red},     % registers
}

\lstdefinestyle{C}
{
  language=C,
  escapechar=@, % include LaTeX code between `@' characters
  keepspaces,   % needed to preserve spacing with lstinline
  basicstyle=\scriptsize\ttfamily\bfseries,
  commentstyle=\color{darkgray},
  stringstyle=\color{cyan},
  showstringspaces=false,
  keywordstyle=[1]\color{blue},    % instructions
  keywordstyle=[2]\color{magenta}, % directives
  keywordstyle=[3]\color{red},     % registers
}

\lstset{%
  inputencoding=utf8,
  extendedchars=true,
  literate=%
  {é}{{\'e}}{1}%
  {è}{{\`e}}{1}%
  {à}{{\`a}}{1}%
  {ç}{{\c{c}}}{1}%
  {œ}{{\oe}}{1}%
  {ù}{{\`u}}{1}%
  {É}{{\'E}}{1}%
  {È}{{\`E}}{1}%
  {À}{{\`A}}{1}%
  {Ç}{{\c{C}}}{1}%
  {Œ}{{\OE}}{1}%
  {Ê}{{\^E}}{1}%
  {ê}{{\^e}}{1}%
  {î}{{\^i}}{1}%
  {ô}{{\^o}}{1}%
  {û}{{\^u}}{1}%
  {ë}{{\¨{e}}}1
  {û}{{\^{u}}}1
  {â}{{\^{a}}}1
  {Â}{{\^{A}}}1
  {Î}{{\^{I}}}1
}

\newcommand{\mycolumns}[5]{
\begin{columns}[#1]
    \column{#2cm}
    #4
    \column{#3cm}
    #5
\end{columns}
}

\newcommand{\hlrose}[1]{\sethlcolor{rosepale}\hl{#1}}
\newcommand{\non}[1]{\ensuremath{\overline{#1}}}
\newcommand{\mytriangle}[6]{
\begin{center}
\begin{tabular}{m{#2cm}m{#1cm}}
    #5 & \includegraphics[width=#1cm]{figs/#4} \\
    \multicolumn{2}{m{#3cm}}{#6} \\
\end{tabular}
\end{center}
}

\title[Assembleur ARM avancé]{Assembleur ARM avancé}
\author[C.~Tibermacine]{\large{Chouki TIBERMACINE}\\
\vspace{1cm}
\small{Chouki.Tibermacine@umontpellier.fr}}
\institute{Polytech Montpellier}
\date{\tiny{IG3 2018-2019}}

\begin{document}

\begin{frame}
\titlepage
\begin{center}
\includegraphics[width=4cm]{figs/polytech.png}
\end{center}
\end{frame}

\begin{frame}
\frametitle{Objectifs du chapitre}

\begin{enumerate}
\item Donner l'intuition de la compilation des langages de haut
  niveau;
\item Comprendre le mécanisme d'appel de fonction;
\item Comprendre l'utilisation de la pile;
\item Étudier la définition et l'utilisation de tableaux.
\end{enumerate}
\end{frame}

\begin{frame}
\frametitle{Haut niveau et bas niveau}

\begin{block}{Traduction assembleur}
\begin{itemize}
\item Des constructions qui peuvent sembler élémentaires dans un
  langage tel que C ou Python (tests, boucles, fonctions) seront
  traduites par de nombreuses instructions assembleur

\item Par ailleurs, des notions d'apparence pourtant simple n'ont pas
  d'équivalent. Par exemple, les variables de Python seront traduites
  en termes d'emplacements mémoire et de registres, une tâche non
  triviale (manipuler des adresses).
\end{itemize}
\end{block}
\end{frame}

\begin{frame}[fragile]
\frametitle{Conditionnelles}

\begin{block}{Code C}
{\small
\begin{lstlisting}[style=C]
if (r1 < r2)
  r3 = r1;
else
  r3 = r2;
\end{lstlisting}}
\end{block}

\begin{block}{Code ARM}
{\small
\begin{lstlisting}[style=arm]
      cmp r1, r2
      blt then        // si r1 < r2 saut vers then
      mov r3, r2      // r3 := r2
      b done          // saut vers done
then: mov r3, r1      // r3 := r1
done:                 // suite du programme
\end{lstlisting}}
\end{block}
\end{frame}

\begin{frame}[fragile]
\frametitle{Boucles}

\begin{block}{Code C}
{\small
\begin{lstlisting}[style=C]
r2 = 0;
while (r1 > 0) {
  r2 = r2 + r1;
  r1 = r1 - 1;
}
\end{lstlisting}}
\end{block}

\begin{block}{Code ARM}
{\small
\begin{lstlisting}[style=arm]
       mov r2, #0     // r2 := 0
while: cmp r1, #0
       ble done       // si r1 <= 0 saut vers done
       add r2, r2, r1 // r2 := r2 + r1
       sub r1, r1, #1  // r1 := r1 - 1
       b while        // retour vers while
done:                 // suite du programme
\end{lstlisting}}
\end{block}
\end{frame}

\begin{frame}[fragile]
\frametitle{Tableaux}
\begin{block}{Déclaration d'un tableau}
  En C : int a[100];
\begin{lstlisting}[style=arm]
	.balign 4
a:	.skip 400
\end{lstlisting}

\end{block}

\begin{block}{Accès à un élément du tableau}
En C : a[3] = 3;
  \begin{lstlisting}[style=arm]
         ldr r1, =a
         mov r2, #3
         str r2, [r1, #+12]
\end{lstlisting}

\end{block}

\begin{block}{Exercice}
  Écrire un programme qui remplit un tableau de 100 éléments avec les
  100 premiers entiers positifs (0$\rightarrow$99) et les affiche ensuite
\end{block}
\end{frame}

\begin{frame}[fragile]
\frametitle{Fonctions}

\begin{block}{Code C}
{\small
\begin{lstlisting}[style=C]
int succ (int x) {
  return x + 1;
}
\end{lstlisting}}
\end{block}

\begin{block}{Code ARM}
{\small
\begin{lstlisting}[style=arm]
/* Fonction succ */
succ:                // r0 contient l'argument x
      add r0, r0, #1 // r0 contient le résultat
      bx lr          // retour à l'appelant (lr : registre qui contient
                     // l'adresse de l'instruction suivant l'appel)

/* Côté appelant */
      mov r0, #10    // On met la valeur de l'argument dans r0
      bl succ        // On appelle la fonction succ
    \end{lstlisting}}
\end{block}
\end{frame}

\begin{frame}[fragile]
\frametitle{Fonctions avec un corps qui modifie le registre lr}

\begin{block}{Code ARM}
{\small
\begin{lstlisting}[style=arm]
      .data
msg:  .string "Nouvelle valeur : %d\n"
      .text
succ: add r1, r0, #1
      ldr r0, =msg  
      bl printf      // Instruction qui modifie lr
      mov r0, r1
      bx lr     

      // Côté appelant :
      mov r0, #10
      bl succ
    \end{lstlisting}}
\end{block}
\only<1>{
  {\color{red}Problème avec l'instruction $bx~lr$ car lr a été modifié par $bl~printf$ (Qu'est-ce qui va se passer à l'exécution ?)}
}
\only<2>{
    {\color{red}Problème avec l'instruction $bx~lr$ car lr a été modifié par $bl~printf$ (boucle infinie, entre les 2 dernières instructions de succ)}
  }

\end{frame}

\begin{frame}[fragile]
\frametitle{Fonctions avec un corps qui modifie le registre lr}

\begin{block}{Solution : enregistrer en mémoire l'@ qui se trouve dans $lr$}
{\small
\begin{lstlisting}[style=arm]
        .data
return: .word 0 // Déclarer une variable return
msg:    .string "Nouvelle valeur : %d\n"
        .text
succ:   ldr r1, =return // Enregistrer dans r1 l'adresse de return
	str lr, [r1] // Enregistrer dans return l'adresse dans lr
        add r1, r0, #1
        ldr r0, =msg  
        bl printf
        mov r0, r1
        ldr r1, =return // Récupérer de return l'adresse enregistrée
	ldr lr, [r1] // Remettre cette adresse dans lr
        bx lr     
      \end{lstlisting}}
\end{block}
\end{frame}

\begin{frame}
\frametitle{Variables locales et appels récursifs}
  \begin{itemize}
  \item Si l'on veut déclarer des variables locales dans une fonction ?
  \item Si l'on veut faire des appels récursifs, comment on va gérer
    les adresses de retour (solution précédente : variable globale,
    écrasée à chaque appel)?
  \end{itemize}

\begin{block}{Solution : utiliser la pile}
  \begin{itemize}
  \item C'est quoi une pile ? c'est un espace en mémoire utilisé lors
    d'un appel de fonction
  \item L'adresse du sommet de la pile est stockée dans le registre
    $sp$ (\textit{stack pointer})
  \item Cet espace est extensible ``à souhait'' par une fonction pour
    stocker des valeurs locales pour ses calculs
  \item Cet espace doit (par convention) être libéré par la fonction
  \end{itemize}
\end{block}
\end{frame}

\begin{frame}[fragile]
  \frametitle{Enregistrement de l'adresse qui se trouve dans lr dans
    la pile}

\begin{block}{Code ARM}
{\small
\begin{lstlisting}[style=arm]                  
foo: sub sp, sp, #8  // sp := sp - 8. On aggrandit la pile de 8 octets
                     // sommet de pile : adresse du dernier octet
                     // d'où la soustraction (sub). La pile grandit
                     // dans le sens des adresses décroissantes.
                     // Pourquoi 8 octets et non 4 ? standard ARM !
     str lr, [sp]    // On enregistre lr en sommet de pile 
     ... // Code de la fonction foo
     ldr lr, [sp]    // On restitue lr depuis la pile 
     add sp, sp, #8  // sp := sp + 8.
                     // pour réduire la taille de la pile de 8 octets
                     // et donc remettre le pointeur de pile à sa valeur
                     // d'origine
     bx lr
\end{lstlisting}}
\end{block}
\end{frame}

\begin{frame}[fragile]
\frametitle{Variables locales}

\begin{block}{Code C}
{\small
\begin{lstlisting}[style=C]
int succ (int x) {
  int y;
  y = x + 1;
  return y;
}
\end{lstlisting}}
\end{block}

\begin{block}{Code ARM}
{\small
\begin{lstlisting}[style=arm]                  
succ: sub sp, sp, #8   // allocation mémoire pour la variable locale
      add r0, r0, 1    // opération
      str r0, [sp]     // enregistrement de la valeur dans la pile
      add sp, sp, #8   // désallocation
      bx lr            // retour
\end{lstlisting}}
\end{block}
\end{frame}
  
\begin{frame}
\frametitle{Convention d'appel}

\begin{block}{Règles}

  La convention d'appel régit la communication entre appelant et
  appelé lors de l'appel de fonction. Pour l'ARM, la convention
  proposée par le fabricant est la suivante:
\begin{footnotesize}
\begin{itemize}
\item les arguments sont passés dans $r0$-$r3$, puis (s'il y en a plus
  de 4) sur la pile;
\item les valeurs de retour sont renvoyées dans $r0$ (et dans
  $r1$-$r3$ si renvoi de plusieurs/grandes valeurs). L'appelant ne
  doit pas faire d'hypothèse sur les valeurs dans ces registres, à
  moins qu'ils soient utilisés pour le passage de paramètres;
\item les registres restants ($r4$-$r11$ et $sp$), peuvent être
  modifiés par l'appelé mais leurs valeurs doivent être rétablies à la
  fin de la fonction (l'appelant fait l'hypothèse que ces registres ne
  sont pas altérés par l'appelé).
\end{itemize}
\end{footnotesize} {\color{red}Exercice : Écrire un programme qui fait
  un appel à $printf$ pour afficher \textit{Hello World}; ensuite
  debugger ce programme avec $gdb$ pour voir l'état du registre r0
  avant et après l'appel à $printf$}
\end{block}
\end{frame}

\begin{frame}[fragile]
  \frametitle{Sujet de TP}
  \begin{block}{Exercices}
\begin{itemize}
\item Écrire une fonction qui permute le contenu de deux variables
  entières de la zone de données au moyen d'une variable locale (en
  utilisant la pile) pour effectuer la permutation;

\item Écrire le code assembleur correspondant au code C suivant:

{\small
\begin{lstlisting}[style=C]
int sqr (int x) {
  return x * x;
}

int sum_sqr (int x, int y) {
  return sqr(x) + sqr(y);
}
\end{lstlisting}}
\end{itemize}
\end{block}
\end{frame}

\begin{frame}[fragile]
\frametitle{Fonctions récursives}

\begin{block}{Code C}
{\small
\begin{lstlisting}[style=C]
int fact (int n) {
  if (n <= 0)
    return 1;
  else
    return n * fact (n - 1);
}
\end{lstlisting}}
\end{block}
\end{frame}

\begin{frame}[fragile]
\frametitle{Fonctions récursives}

\begin{block}{Code ARM}
{\footnotesize
\begin{lstlisting}[style=arm]
fact:	sub sp, sp, #8
        str lr, [sp]   // enregistrer lr en sommet de pile
	cmp r0, #0
	ble fact0      // si r0 <= 0 saut vers fact0
	sub r0, #1     // r0 := r0 - 1
	mov r1, r0     // enregistrer r0 dans r1
	b fact         // appel récursif
endf:  	
	mul r0, r1     // r0 := r0 * r1 (r1 : résultat précédent)
	ldr lr, [sp]
	add sp, sp, #8
	bx lr          // retour à l'appelant
	
fact0:	mov r0, #1     // valeur de retour = 0
	b endf
	
\end{lstlisting}}

Ce code est erroné. Où est l'erreur?
\end{block}
\end{frame}

\begin{frame}
\frametitle{Fonctions récursives}

\begin{block}{Pourquoi le code est-il erroné?}
\begin{itemize}
\item Parce que fact est récursive, plusieurs appels à la même
  fonction qui s'accumulent (appels actifs en même temps).

\item L'appel récursif modifie $r0$ et $r1$; en fait, il en détruit le
  contenu. Or, celui-ci est utilisé après l'appel.

\item Il faut donc sauvegarder le contenu de $r0$ avant l'appel, puis
  le restaurer après l'appel (comme pour $lr$).
\end{itemize}
\end{block}
\end{frame}

\begin{frame}[fragile]
\frametitle{Fonctions récursives}
\vspace{-1cm}
\begin{block}{Code ARM}
  \vspace{-.2cm}
{\small
\begin{lstlisting}[style=arm]
fact:	sub sp, sp, #8
	str lr, [sp]
	cmp r0, #0
	ble fact0
	sub sp, sp, #8 
	str r0, [sp]   // Enregistrer r0 dans la pile (empiler r0)
	sub r0, #1     
	bl fact        // Appel récursif avec r0-1
endf:	ldr r1, [sp]   // Récupérer dans r1 la valeur empilée
	add sp, sp, #8
	mul r0, r1     // Multiplier r0 par r1 (résultat dans r0)
	ldr lr, [sp]
	add sp, sp, #8
	bx lr          // Retourner vers endf ou vers l'appelant de fact
fact0:	mov r0, #1     
	sub sp, sp, #8
	str r0, [sp]   // Empiler la valeur 1 pour fact(0)
	b endf
      \end{lstlisting}}
    Dérouler l'exécution d'un programme avec un appel : $fact(3)$
\end{block}
\end{frame}

\begin{frame}[fragile]
\frametitle{Sujet de TP -suite-}

\begin{block}{Exercices}
\begin{itemize}
\item Écrire une fonction qui effectue récursivement la somme des $n$
  premiers entiers, où $n$ est un entier passé en argument;

\item Écrire la fonction de Fibonacci:

$$fib(n)=
\left\{
\begin{array}{l}
1\mbox{, si }n=0,1\\
fib(n-1)+fib(n-2)\mbox{, sinon}
\end{array}
\right.$$
\end{itemize}
\end{block}
\end{frame}

% \begin{frame}[fragile]
% \frametitle{Plus simplement}

% \begin{block}{Appels de fonctions dans les fonctions}
% \begin{itemize}
% \item Le registre $lr$ sera écrasé par un sous-appel;
% \item Il faut donc le sauvegarder systématiquement;
% \item Même s'il n'y a pas de sous-appel;
% \item Cela permet d'éviter de devoir y penser.
% \end{itemize}
% \end{block}

% \begin{block}{Toute fonction sera parenthésée par}
% {\small
% \begin{lstlisting}[style=arm]
% func: sub sp, sp, #8
%       str lr, [sp]
%       ...
%       ldr lr, [sp]
%       add sp, sp, #8
%       bx lr
% \end{lstlisting}}
% \end{block}
% \end{frame}

\end{document}
